\chapter[Zanurzenia i zwartość Sobolewa]{Zanurzenia Sobolewa, nierówność Morreya, zwartość i ślady}
\label{chap:pde-embeddings}

Przestrzeń $W^{1,p}(\Omega)$ została zdefiniowana przez całkowalność
funkcji i jej słabych pochodnych. Nie wynika z tej definicji wprost,
czy funkcja jest ciągła, czy można ją przybliżać jednostajnie
ani czy ciąg ograniczony w $W^{1,p}$ ma podciąg zbieżny w $L^p$.
Odpowiedzi zależą przede wszystkim od relacji $p$ do wymiaru $d$.
Będziemy odróżniać \emph{zanurzenie ciągłe}, oznaczane
$X\hookrightarrow Y$, od \emph{zanurzenia zwartego}, oznaczanego
$X\hookrightarrow\hookrightarrow Y$. Pierwsze znaczy, że
$\|u\|_Y\le C\|u\|_X$; drugie, że każdy ciąg ograniczony w $X$
ma podciąg zbieżny w normie $Y$.

W rozdziale $d\in\N$, $d\ge1$, i $1\le p\le\infty$.
Przyjmujemy, że $\Omega$ jest niepustym otwartym podzbiorem
$\R^d$, a funkcje są rzeczywiste, o ile nie podano inaczej.
Dla gradientu używamy punktowej normy euklidesowej.
Norma ta jest równoważna normie składowych używanej
w rozdziale~\ref{chap:pde-sobolev}: dla $p<\infty$ wynika
to z równoważności norm $\ell^2$ i $\ell^p$ na $\R^d$,
po podniesieniu do potęgi $p$ i scałkowaniu, a dla
$p=\infty$ z nierówności
$\max_j|v_j|\le |v|\le\sqrt d\max_j|v_j|$.
Nie zakładamy automatycznie,
że $\Omega$ ma regularny brzeg. Kiedy korzystamy z obszaru Lipschitzowskiego, podajemy to wyraźnie. Całki na $\R^d$ i na
$\Omega$ są całkami względem $\lambda_d$; na brzegu pojawia się
miara powierzchniowa $\mathcal H^{d-1}$.

\section{Nierówność geometryczna potrzebna do dowodu Sobolewa}
\label{sec:pde10-LW}

Zaczniemy od rachunkowego lematu, który pozwala udowodnić
nierówność Sobolewa bez odwoływania się do twierdzenia
izoperymetrycznego ani do teorii potencjałów Riesza.
Dla $x=(x_1,\ldots,x_d)$ przez $\widehat x_j$ oznaczamy
wektor powstały przez opuszczenie współrzędnej $x_j$.

\begin{lemat}[Nierówność Loomisa--Whitneya w postaci całkowej]
\label{lem:pde10-LW}
Niech $d\ge2$ i $F_j:\R^{d-1}\to[0,\infty]$ będą mierzalne,
$j=1,\ldots,d$. Wówczas
\[
 \int_{\R^d}\prod_{j=1}^d F_j(\widehat x_j)^{1/(d-1)}
 \,\mathrm d\lambda_d(x)
 \le \prod_{j=1}^d
 \left(\int_{\R^{d-1}}F_j\,\mathrm d\lambda_{d-1}\right)^{1/(d-1)}.
\]
Jeżeli pewna całka po prawej stronie jest zerowa, to
lewa strona również jest zerowa, nawet gdy inne całki
są nieskończone. W tym przypadku przyjmujemy prawą
stronę równą zeru. Jeżeli wszystkie całki są dodatnie,
a przynajmniej jedna nieskończona, prawa strona jest
nieskończona. W iloczynie podcałkowym przyjmujemy wartość
zero, gdy którykolwiek czynnik jest zerowy, również
w obecności czynników nieskończonych; odpowiada to
granicy iloczynów nieujemnych obcięć.
\end{lemat}
\begin{proof}
Wystarczy najpierw rozważyć nieujemne funkcje ograniczone
o zwartych nośnikach. Dla $d=2$ nierówność jest równością:
$F_1$ zależy tylko od $x_2$, a $F_2$ tylko od $x_1$.
Z twierdzenia Tonellego otrzymujemy
\[
 \int_{\R^2}F_1(x_2)F_2(x_1)\,\mathrm d\lambda_2(x)
 =\left(\int_\R F_1\,\mathrm d\lambda_1\right)
  \left(\int_\R F_2\,\mathrm d\lambda_1\right).
\]
W rozważanym przypadku ograniczonych funkcji o zwartych
nośnikach wszystkie całki są skończone.

Załóżmy indukcyjnie tezę w wymiarze $d-1$, gdzie $d\ge3$.
Całkujemy w lewej stronie po $x_d$. Czynnik
$F_d(x_1,\ldots,x_{d-1})$ nie zależy od $x_d$.
Nierówność H\"oldera zastosowana do pozostałych $d-1$
czynników, każdy z wykładnikiem $d-1$, daje
\[
 \int_{\R}\prod_{j<d}F_j(\widehat x_j)^{1/(d-1)}
 \,\mathrm d\lambda_1(x_d)
 \le\prod_{j<d}G_j(\widehat x^{\prime}_j)^{1/(d-1)},
\]
gdzie $x'=(x_1,\ldots,x_{d-1})$, a $G_j$ powstaje przez
scałkowanie $F_j$ względem $x_d$; nie zależy więc od $x_j$.
Oznaczmy $H(x')=\prod_{j<d}G_j(\widehat x'_j)^{1/(d-2)}$.
Zapiszmy dla skrócenia wzorów
$I_d=\int_{\R^{d-1}}F_d(x')\,\mathrm d\lambda_{d-1}(x')$,
$J=\int_{\R^{d-1}}H(x')\,\mathrm d\lambda_{d-1}(x')$
oraz, dla $j<d$,
$I_j=\int_{\R^{d-2}}G_j(z)\,\mathrm d\lambda_{d-2}(z)$.
Po użyciu H\"oldera na $\R^{d-1}$ z wykładnikami
$d-1$ oraz $(d-1)/(d-2)$ otrzymujemy
\begin{align*}
 &\int_{\R^{d-1}}F_d(x')^{1/(d-1)}
 \prod_{j<d}G_j(\widehat x'_j)^{1/(d-1)}
 \,\mathrm d\lambda_{d-1}(x')\\
 &\qquad\le I_d^{1/(d-1)}J^{(d-2)/(d-1)}
 \le I_d^{1/(d-1)}\prod_{j<d} I_j^{1/(d-1)}.
\end{align*}
Druga nierówność jest dokładnie założeniem indukcyjnym.
Na mocy twierdzenia Fubiniego
$I_j=\int_{\R^{d-1}}F_j(z)\,\mathrm d\lambda_{d-1}(z)$,
co kończy krok indukcyjny. Dla dowolnych nieujemnych $F_j$
kładziemy
$F_{j,n}=\min(F_j,n)\chi_{[-n,n]^{d-1}}$.
Funkcje te są ograniczone, mają nośniki zawarte w zbiorach
zwartych i rosną punktowo do $F_j$. Iloczyny
$\prod_j F_{j,n}(\widehat x_j)^{1/(d-1)}$ również rosną
do iloczynu z tezy zgodnie z podaną konwencją.
Twierdzenie o zbieżności monotonicznej stosujemy do lewej
strony i do każdej całki po prawej stronie.
Jeżeli $\int F_j=0$ dla pewnego $j$, to wszystkie całki
$\int F_{j,n}$ są zerowe, więc lewa strona jest zerowa
w granicy. W przeciwnym razie wszystkie granice całek
są dodatnie; iloczyn prawych stron zbiega do ich iloczynu,
również gdy jedna z granic jest nieskończona.
To dowodzi nierówności w pełnej ogólności.
\end{proof}

\begin{twierdzenie}[Nierówność Sobolewa dla $p=1$]
\label{thm:pde10-Sob1}
Dla $d\ge2$ i każdej $u\in C_c^1(\R^d)$ zachodzi
\[
 \|u\|_{L^{d/(d-1)}(\R^d)}
 \le\prod_{j=1}^d\|\partial_j u\|_{L^1(\R^d)}^{1/d}
 \le \frac1d\sum_{j=1}^d\|\partial_j u\|_{L^1(\R^d)}.
\]
W szczególności $\|u\|_{L^{d/(d-1)}}\le C_d\|\nabla u\|_{L^1}$.
\end{twierdzenie}
\begin{proof}
Dla ustalonego $j$ i każdego $x\in\R^d$ funkcja jednej
zmiennej $t\mapsto u(x_1,\ldots,t,\ldots,x_d)$ jest
klasy $C^1$ i znika dla dostatecznie ujemnych $t$ dzięki
zwartości nośnika $u$. Z podstawowego twierdzenia
rachunku całkowego
\[
 u(x)=\int_{-\infty}^{x_j}
 \partial_j u(x_1,\ldots,t,\ldots,x_d)\,\mathrm d\lambda_1(t).
\]
Biorąc moduł i rozszerzając przedział całkowania
nieujemnej funkcji, uzyskujemy
\[
 |u(x)|\le A_j(\widehat x_j),\qquad
 A_j(\widehat x_j)=\int_{\R}
 |\partial_j u(x_1,\ldots,t,\ldots,x_d)|\,\mathrm d\lambda_1(t).
\]
Mnożąc $d$ nierówności i podnosząc do potęgi $1/(d-1)$,
dostajemy $|u(x)|^{d/(d-1)}\le\prod_j
A_j(\widehat x_j)^{1/(d-1)}$. Lemat~\ref{lem:pde10-LW}
daje zatem
\[
 \|u\|_{d/(d-1)}^{d/(d-1)}
 \le \prod_{j=1}^d\|\partial_j u\|_1^{1/(d-1)}.
\]
Po podniesieniu do potęgi $(d-1)/d$ otrzymujemy pierwszą
nierówność. Druga jest nierównością pomiędzy średnią
geometryczną i arytmetyczną. Ponadto z nierówności
Cauchy'ego--Schwarza punktowo
$\sum_j|\partial_j u|\le\sqrt d\,|\nabla u|$.
Całkując, możemy więc przyjąć $C_d=1/\sqrt d$.
Wszystkie całki są skończone
z uwagi na zwarty nośnik $u$.
\end{proof}

\section{Ostre wykładniki zanurzenia Sobolewa}
\label{sec:pde10-Sobolev}

\begin{twierdzenie}[Nierówność Sobolewa na całej przestrzeni]
\label{thm:pde10-Sobp}
Niech $d\ge2$ i $1\le p<d$. Przyjmijmy
$p^*=dp/(d-p)$. Istnieje stała $C(d,p)>0$ taka, że
\[
 \|u\|_{L^{p^*}(\R^d)}
 \le C(d,p)\|\nabla u\|_{L^p(\R^d)}
 \quad\text{dla każdego }u\in C_c^\infty(\R^d).
\]
W konsekwencji identyczne oszacowanie obowiązuje
(z ewentualnie inną stałą) dla $u\in W^{1,p}(\R^d)$.
\end{twierdzenie}
\begin{proof}
Dla $p=1$ jest to twierdzenie~\ref{thm:pde10-Sob1}.
Niech $1<p<d$ i $u\in C_c^\infty(\R^d)$. Wybierzmy
$\gamma=p(d-1)/(d-p)>1$. Stosujemy twierdzenie
\ref{thm:pde10-Sob1} do funkcji
$v=|u|^\gamma$. Ponieważ $\gamma>1$, funkcja
$t\mapsto |t|^\gamma$ jest klasy $C^1$ na $\R$:
dla $t\ne0$ jej pochodna wynosi
$\gamma|t|^{\gamma-1}\operatorname{sgn}(t)$,
a w zerze iloraz różnicowy dąży do zera. Pochodna
jest ciągła także w zerze. Zatem $v\in C_c^1(\R^d)$,
gdyż $u$ jest gładka i $v=0$ poza nośnikiem $u$.
Jest to dokładnie klasa funkcji dopuszczalna w twierdzeniu
dla $p=1$.
Reguła łańcuchowa oraz nierówność H\"oldera dają
\begin{align*}
 \|u\|_{L^{\gamma d/(d-1)}}^\gamma
 &\le C_d\gamma\int_{\R^d}|u|^{\gamma-1}|\nabla u|
 \,\mathrm d\lambda_d\\
 &\le C_d\gamma
 \|u\|_{L^{(\gamma-1)p'}}^{\gamma-1}
 \|\nabla u\|_{L^p},\qquad p'=\frac p{p-1}.
\end{align*}
Z wyboru $\gamma$ wynika
\[
 \gamma\frac d{d-1}=\frac{dp}{d-p}=p^*,
 \qquad
 \gamma-1=\frac{d(p-1)}{d-p},
 \qquad
 (\gamma-1)p'=\frac{dp}{d-p}=p^*.
\]
Jeśli $u\not\equiv0$, dzielimy obie strony przez
$\|u\|_{p^*}^{\gamma-1}$; przypadek $u=0$ jest oczywisty.
To daje nierówność dla funkcji gładkich.

Jeśli $u\in W^{1,p}(\R^d)$, wybieramy
$u_n\in C_c^\infty(\R^d)$ zbieżne do $u$ w $W^{1,p}$.
Istnienie takiego ciągu wynika z gęstości
$C_c^\infty(\R^d)$ w $W^{1,p}(\R^d)$ udowodnionej
w sekcji~\ref{sec:pde9-density}; tu $p<\infty$.
Stosujemy oszacowanie do różnicy $u_n-u_m$:
\[
 \|u_n-u_m\|_{p^*}
 \le C(d,p)\|\nabla u_n-\nabla u_m\|_p\longrightarrow0.
\]
Zupełność $L^{p^*}$ daje granicę $v$.
Dla dowolnego $t>0$ nierówność trójkąta i nierówność
Czebyszewa dają
\[
 \lambda_d(\{|u-v|>t\})
 \le (2/t)^p\|u-u_n\|_p^p
     +(2/t)^{p^*}\|u_n-v\|_{p^*}^{p^*}.
\]
Prawa strona dąży do zera, zatem, biorąc $t=1/m$,
$m\in\N$, otrzymujemy $u=v$ prawie wszędzie.
Ciągłość normy po obu stronach nierówności dla $u_n$
daje tezę dla $u$ z tą samą stałą $C(d,p)$.
\end{proof}

\begin{uwaga}[Dlaczego wykładnik jest krytyczny?]
\label{rem:pde10-critical}
Niech $d\ge2$, $1\le p<d$ oraz $1\le q\le\infty$.
Dla $\varphi\in C_c^\infty(B_1(0))$, $\varphi\ne0$, oraz
$x_0\in\R^d$ niech
$u_\varepsilon(x)=\varepsilon^{1-d/p}
\varphi((x-x_0)/\varepsilon)$. Wówczas
\[
 \|\nabla u_\varepsilon\|_p=\|\nabla\varphi\|_p,
 \qquad \|u_\varepsilon\|_q
 =\varepsilon^{1-d/p+d/q}\|\varphi\|_q.
\]
Wzory wynikają z podstawienia $x=x_0+\varepsilon y$,
reguły łańcuchowej oraz (dla $q=\infty$) zachowania
zbiorów miary zero, jak w sekcji~\ref{sec:pde9-scaling}.
W szczególności $\|u_\varepsilon\|_p=\varepsilon\|\varphi\|_p$,
więc normy $W^{1,p}$ są jednostajnie ograniczone dla
$0<\varepsilon\le1$. Dla $q=p^*$ drugi wykładnik
jest równy zeru; dla $q>p^*$ jest ujemny, a wówczas
$\|u_\varepsilon\|_q\to\infty$ przy $\varepsilon\downarrow0$.
Nie może więc istnieć
oszacowanie $\|u\|_q\le C\|u\|_{W^{1,p}}$ dla
wszystkich takich funkcji, nawet na dowolnym obszarze zawierającym kulę wokół $x_0$. Ten sam przykład pokaże
brak zwartości w wykładniku krytycznym.
\end{uwaga}

\begin{definicja}[Obszar Lipschitzowski]
\label{def:pde10-Lipschitz}
Ograniczony otwarty zbiór $\Omega\subset\R^d$ nazywamy
\emph{obszarem Lipschitzowskim}, jeżeli jest spójny, a jego
brzeg można pokryć skończoną liczbą otoczeń, w których
po odpowiednim sztywnym ruchu współrzędnych ma on postać
$\{(x',x_d):x_d=\phi(x')\}$ z $\phi$ Lipschitzowską,
zaś $\Omega$ leży lokalnie po jednej stronie tego wykresu.
Przyjmujemy, że otoczenia są mniejsze niż wykresowe
cylindry, aby dostępne były dodatnie odległości między
nośnikami funkcji odcinających a ich bocznymi ścianami.
\end{definicja}

\begin{lemat}[Przedłużanie pierwszego rzędu]
\label{lem:pde10-extension}
Niech $\Omega\subset\R^d$ będzie obszarem ograniczonym Lipschitzowskim i $1\le p\le\infty$. Istnieje liniowy operator
$E:W^{1,p}(\Omega)\to W^{1,p}(\R^d)$ oraz stała
$C=C(\Omega,p)>0$ takie, że $(Eu)|_\Omega=u$ i
$\|Eu\|_{W^{1,p}(\R^d)}\le C\|u\|_{W^{1,p}(\Omega)}$.
Można dodatkowo wybrać operator o obrazie w funkcjach
wspartych w jednej ustalonej kuli, niezależnej od $u$.
\end{lemat}
\begin{proof}
\emph{Krok 1. Prostowanie brzegu.}
Zmniejszając otoczenie wykresu, możemy przyjąć
$U=\{(x',x_d):x'\in D,\ |x_d-\phi(x')|<h\}$,
gdzie $D$ jest kulą, $h>0$, a $\Omega\cap U$
jest dokładnie częścią $x_d>\phi(x')$ tego otoczenia
(w razie potrzeby odwracamy ostatnią oś).
Niech $Q=D\times(-h,h)$, $Q_+=D\times(0,h)$ oraz
\[
 \Phi(x',t)=(x',t+\phi(x')), \qquad
 \Psi(x',x_d)=(x',x_d-\phi(x')).
\]
Są to wzajemnie odwrotne odwzorowania dwulipschitzowskie,
co wynika z nierówności trójkąta i własności Lipschitza
funkcji $\phi$. Mamy $\Phi(Q_+)=U\cap\Omega$.
Dla nieujemnej funkcji borelowskiej $a$ z Tonellego
i podstawienia w ostatniej współrzędnej otrzymujemy
\[
 \int_Q a\circ\Phi\,\mathrm d\lambda_d
 =\int_D\int_{\phi(x')-h}^{\phi(x')+h}a(x',s)
 \,\mathrm d\lambda_1(s)\,\mathrm d\lambda_{d-1}(x')
 =\int_U a\,\mathrm d\lambda_d.
\]
Analogicznie dla $Q_+$ i $U\cap\Omega$.
Wzór dla funkcji charakterystycznych borelowskich
zbiorów zerowych daje zachowanie zbiorów miary zero
w obu kierunkach; przez wybór reprezentantów borelowskich
rozszerzamy wzór na funkcje mierzalne w sensie Lebesgue'a.
Składanie zachowuje normy $L^p$, także dla $p=\infty$.

Dla gładkiej $v$ na $U\cap\Omega$ ograniczenia $\phi$
do prostych równoległych do osi są lipschitzowskie,
a więc absolutnie ciągłe na zwartych odcinkach.
Z twierdzenia~\ref{thm:pde9-ac} i Fubiniego pochodne
$\partial_i\phi$ istnieją prawie wszędzie na $D$
i są ograniczone co do modułu przez $L=\operatorname{Lip}\phi$.
Są mierzalne jako granice ilorazów różnicowych
(poza zbiorem istnienia przyjmujemy zero).
Na zwartych pododcinkach prostych w $Q_+$ złożenie
$v\circ\Phi$ jest absolutnie ciągłe, ponieważ krzywa
wewnętrzna jest lipschitzowska, a $v$ lokalnie
lipschitzowska na jej zwartym obrazie.
Jednowymiarowa reguła łańcuchowa daje prawie wszędzie
\[
 \partial_t(v\circ\Phi)=(\partial_d v)\circ\Phi,\qquad
 \partial_i(v\circ\Phi)=(\partial_i v)\circ\Phi+
    ((\partial_d v)\circ\Phi)\partial_i\phi\quad(i<d).
\]
Są to również słabe pochodne: całkujemy przez części
na odcinkach przeciw testowi o zwartym nośniku,
a następnie stosujemy Fubiniego. Test znika przy końcach
przekrojów; całki są bezwzględnie zbieżne dzięki
lokalnej ograniczoności czynników na nośniku testu.
Nie korzystamy z twierdzenia o przekrojach ogólnych
funkcji Sobolewa. Z zachowania miary wynika
\[
 \|\partial_t(v\circ\Phi)\|_p=\|\partial_d v\|_p,\qquad
 \|\partial_i(v\circ\Phi)\|_p
 \le\|\partial_i v\|_p+L\|\partial_d v\|_p.
\]
Normy po lewej liczymy na $Q_+$, po prawej na $U\cap\Omega$.
Dla $\Psi$ rachunek jest analogiczny ze znakiem minus
przy $\partial_i\phi$; działa też na pełnych otoczeniach.
Jeśli przekrój ma wiele składowych, całkujemy przez części
na każdej z nich, a całki sumujemy przez bezwzględną
całkowalność.

Dla $p<\infty$ i $v\in W^{1,p}$ wybieramy gładkie
$v_n$ zbieżne do $v$ w tej normie na mocy twierdzenia
~\ref{thm:pde9-MeyersSerrin}. Oszacowania dla różnic
$v_n-v_m$ dają ciąg Cauchy'ego złożeń w $W^{1,p}$.
Zupełność i zachowanie normy $L^p$ identyfikują granicę
jako złożenie z $v$; wzory pochodnych przechodzą
do granicy w $L^p$, gdyż $\partial_i\phi$ jest ograniczona.
Dotyczy to obu odwzorowań oraz pełnych i jednostronnych
otoczeń. Dla $p=\infty$ stosujemy przypadek $p=1$,
ponieważ otoczenia mają skończoną miarę. Prawe strony
wzorów należą do $L^\infty$, co daje te same oszacowania
bez aproksymacji w normie $W^{1,\infty}$ i bez korzystania
z nierówności Morreya.
Dla $d=1$ pomijamy zmienne styczne; zmiana współrzędnych
jest przesunięciem lub odbiciem.

\emph{Krok 2. Przedłużenie lokalne.}
Dla $\chi\in C_c^\infty(U)$ połóżmy
$v=(\chi u)|_{U\cap\Omega}$, $w=v\circ\Phi$ na $Q_+$.
Reguła iloczynu (twierdzenie~\ref{thm:pde9-products})
i krok 1 dają oszacowanie normy $w$ przez stałą razy
$\|u\|_{W^{1,p}(\Omega)}$.
Funkcja $w$ i jej słabe pochodne znikają prawie wszędzie
poza $A=\Psi(\operatorname{supp}\chi)\Subset Q$
(dla pochodnych wynika to z lokalności).
Przedłużamy $w$ przez zero na półprzestrzeń
$P=\R^{d-1}\times(0,\infty)$, oznaczając wynik przez $W$.
Wybierzmy $\zeta\in C_c^\infty(Q)$ równą $1$ w otoczeniu
$A$. Dla $\varphi\in C_c^\infty(P)$ funkcja
$(\zeta\varphi)|_{Q_+}$ jest testem na $Q_+$,
gdyż $\varphi$ znika blisko $t=0$, a $\zeta$ blisko
sztucznych ścian. Z zanikania $w$ i jej pochodnych
poza $A$ wynika
\[
 \begin{aligned}
 -\int_P W\partial_j\varphi\,\mathrm d\lambda_d
 &=-\int_{Q_+}w\partial_j(\zeta\varphi)\,\mathrm d\lambda_d\\
 &=\int_{Q_+}(\partial_jw)\zeta\varphi\,\mathrm d\lambda_d
 =\int_{Q_+}(\partial_jw)\varphi\,\mathrm d\lambda_d.
 \end{aligned}
\]
Pochodne $W$ są zatem przedłużeniami przez zero
pochodnych $w$, z zachowaniem normy.

Odbijamy $W$ parzyście względem $t=0$ na mocy
twierdzenia~\ref{thm:pde9-reflection}. Dla $p<\infty$
norma operatora wynosi $2^{1/p}$. Dla $p=\infty$
funkcja $W$ i jej pochodne są ograniczone i znikają
poza zbiorem ograniczonym, więc $W\in W^{1,1}(P)$.
Stosując twierdzenie dla $p=1$, dostajemy wzory
na odbite pochodne (parzyste styczne i nieparzystą
normalną), które należą do $L^\infty$; odbicie
zachowuje normę maksimum składowych $W^{1,\infty}$.
Odbita funkcja ma nośnik zawarty w
$A\cup\{(x',-t):(x',t)\in A\}\Subset Q$.
Ograniczamy ją do $Q$ i składamy z $\Psi$.
Otrzymujemy funkcję z $W^{1,p}(U)$ o zwartym nośniku
w $U$, którą przedłużamy przez zero za pomocą
lematu~\ref{lem:pde9-compact-product} (z odcięciem
równym $1$ w otoczeniu tego nośnika).
Otrzymany operator $E_\chi$ jest liniowy i ograniczony,
a jego restrykcja do $\Omega$ jest równa $\chi u$.

\emph{Krok 3. Sklejenie.}
Ze zwartości brzegu wybieramy skończenie wiele mniejszych
otoczeń brzegowych o domknięciach zawartych w odpowiednich
$U_j$, $j=1,\ldots,N$. Pozostała część $\overline\Omega$
jest zwartym podzbiorem $\Omega$; obejmujemy ją otwartym
$U_0\Subset\Omega$. Wybieramy nieujemne gładkie odcięcia
$b_j$ o zwartych nośnikach w $U_j$, równe $1$ na
odpowiednich mniejszych zbiorach (konstrukcja ze splotem
z dowodu twierdzenia~\ref{thm:pde9-MeyersSerrin}).
Suma $S=\sum_jb_j$ jest dodatnia na otoczeniu $V$
zbioru $\overline\Omega$. Wybieramy $\eta\in C_c^\infty(V)$
równą $1$ w otoczeniu $\overline\Omega$ i definiujemy
$\chi_j=\eta b_j/S$ na $V$, zero poza $V$.
Są globalnie gładkie, mają zwarte nośniki w $U_j$
i ich suma jest równa $1$ na $\Omega$.
Część $\chi_0u$ przedłużamy przez zero z lematu
~\ref{lem:pde9-compact-product}, pozostałe przez $E_{\chi_j}$,
uwzględniając sztywne ruchy zachowujące miarę
i mające stałe macierze pochodnych.
Suma $Eu$ jest liniowa, a na $\Omega$ równa $u$.
Nierówność trójkąta daje
\[
 \|Eu\|_{W^{1,p}(\R^d)}
 \le\left(\sum_{j=0}^NC_j\right)\|u\|_{W^{1,p}(\Omega)}.
\]
Nośniki składników leżą w ustalonych zbiorach zwartych,
których skończona suma mieści się w jednej kuli.
\end{proof}

\begin{wniosek}[Gęstość funkcji gładkich aż do brzegu]
\label{cor:pde10-density-boundary}
Niech $\Omega\subset\R^d$ będzie ograniczonym obszarem
Lipschitzowskim, $1\le p<\infty$. Wówczas zbiór
$\{\varphi|_\Omega:\varphi\in C_c^\infty(\R^d)\}$
jest gęsty w $W^{1,p}(\Omega)$. W szczególności jest
gęsty zbiór funkcji gładkich w otoczeniu $\overline\Omega$.
\end{wniosek}
\begin{proof}
Dla $u\in W^{1,p}(\Omega)$ przedłużenie
$v=Eu\in W^{1,p}(\R^d)$ istnieje na mocy
lematu~\ref{lem:pde10-extension}. Wybierzmy funkcje
$\chi_R$ równe $1$ w $B_R$, o nośniku w $B_{2R}$
i $\|\nabla\chi_R\|_\infty\le C/R$.
Wtedy $\chi_Rv\to v$ w $W^{1,p}(\R^d)$:
dla składników bez pochodnych korzystamy ze zbieżności
ogonów całek, a dla gradientów dodatkowy składnik
$v\nabla\chi_R$ ma normę najwyżej $C\|v\|_p/R$.
Dla ustalonego $R$ regularizujemy $\chi_Rv$ splotem
z $\rho_\varepsilon$. Ponieważ $p<\infty$,
$\rho_\varepsilon*(\chi_Rv)\to\chi_Rv$ w $W^{1,p}(\R^d)$.
Funkcje te należą do $C_c^\infty(\R^d)$.
Ograniczenie do $\Omega$ nie zwiększa normy, więc
odpowiedni ciąg restrykcji zbiega do $u$ w $W^{1,p}(\Omega)$.
\end{proof}

\begin{twierdzenie}[Zanurzenia na obszarze ograniczonym Lipschitzowskim]
\label{thm:pde10-domain-embed}
Niech $d\ge2$ i $\Omega\subset\R^d$ będzie obszarem ograniczonym Lipschitzowskim. Wówczas zachodzą zanurzenia ciągłe:
\begin{enumerate}[label=\textnormal{(\roman*)}]
\item gdy $1\le p<d$, $W^{1,p}(\Omega)\hookrightarrow
 L^q(\Omega)$ dla $1\le q\le p^*=dp/(d-p)$;
\item gdy $p=d$, $W^{1,d}(\Omega)\hookrightarrow L^q(\Omega)$
dla każdego skończonego $q\ge1$;
\item gdy $p>d$, $W^{1,p}(\Omega)\hookrightarrow
 C^{0,\alpha}(\overline\Omega)$ dla $\alpha=1-d/p$.
Dla $p=\infty$ przyjmujemy $\alpha=1$.
\end{enumerate}
W części (iii) przestrzeń $C^{0,\alpha}(\overline\Omega)$
oznacza klasy z reprezentantem ciągłym o skończonej normie
$\|u\|_\infty+[u]_{C^{0,\alpha}}$.
\end{twierdzenie}
\begin{proof}
Część (i) wynika z lematu~\ref{lem:pde10-extension},
twierdzenia~\ref{thm:pde10-Sobp} i ograniczoności miary
$\Omega$: dla $q\le p^*$ nierówność H\"oldera daje
$\|u\|_q\le|\Omega|^{1/q-1/p^*}\|u\|_{p^*}$.

Gdy $p=d$, ustalmy dowolne skończone $q>d$.
Możemy dobrać $p_0\in(1,d)$ tak, że
$p_0^*=dp_0/(d-p_0)\ge q$. Przedłużenie $Eu$ ma
zwarty nośnik zawarty w ustalonej kuli $B_R$.
Ponieważ $p_0<d$, H\"older na $B_R$ daje
$\|Eu\|_{W^{1,p_0}}\le C\|Eu\|_{W^{1,d}}$.
Twierdzenie~\ref{thm:pde10-Sobp} daje
$\|Eu\|_{p_0^*}\le C\|\nabla Eu\|_{p_0}$.
Po ograniczeniu do $\Omega$ i zastosowaniu Höldera
na tym zbiorze skończonej miary otrzymujemy
$\|u\|_q\le C_q\|u\|_{W^{1,d}(\Omega)}$.
Dla $q\le d$ wystarcza H\"older na $\Omega$.
Część (iii) wynika z nierówności Morreya,
której pełny dowód podajemy w następnej sekcji.
\end{proof}

\section{Nierówność Morreya: dowód z całkowania po odcinkach}
\label{sec:pde10-Morrey}

\begin{lemat}[Oszacowanie różnicy punktu i średniej]
\label{lem:pde10-point-average}
Niech $d\ge2$, $B=B_R(a)\subset\R^d$, a $f\in C^1(\overline B)$.
Dla każdego $x\in B$ zachodzi
\[
 |f(x)-f_B|\le C_d\int_B
 \frac{|\nabla f(z)|}{|x-z|^{d-1}}\,\mathrm d\lambda_d(z),
 \quad f_B=\frac{1}{|B|}\int_Bf\,\mathrm d\lambda_d.
\]
Stała $C_d$ zależy tylko od wymiaru.
\end{lemat}
\begin{proof}
Z wypukłości kuli wynika, że odcinek od $x$ do każdego
$y\in B$ leży w $B$. Stosujemy podstawowe twierdzenie
rachunku całkowego do funkcji $t\mapsto f(x+t(y-x))$.
Z reguły łańcuchowej i nierówności Cauchy'ego--Schwarza
otrzymujemy
$|f(x)-f(y)|\le|x-y|\int_0^1
|\nabla f(x+t(y-x))|\,\mathrm d\lambda_1(t)$.
Ponieważ
$|f(x)-f_B|\le |B|^{-1}\int_B|f(x)-f(y)|\,\mathrm d\lambda_d(y)$,
uśredniamy prawą stronę i zamieniamy całki nieujemne
na mocy Tonellego. Przy ustalonym $t\in(0,1]$ podstawiamy
$z=x+t(y-x)$, zatem $\mathrm d\lambda_d(y)=
t^{-d}\,\mathrm d\lambda_d(z)$ i
$|x-y|=|x-z|/t$. Obrazem $B$ jest kula
$x+t(B-x)\subset B$; każdy jej punkt spełnia
$|x-z|\le 2Rt$. Dla ustalonego $z\ne x$ możemy więc
powiększyć zbiór dopuszczalnych $t$ do przedziału
$[|x-z|/(2R),1]$. Ponownie stosując Tonellego,
otrzymujemy
\begin{align*}
 |f(x)-f_B|
 &\le\frac1{|B|}\int_B|\nabla f(z)|\,|x-z|
 \int_{|x-z|/(2R)}^1 t^{-d-1}\,\mathrm d\lambda_1(t)
 \,\mathrm d\lambda_d(z)\\
 &\le \frac{(2R)^d}{d|B|}
 \int_B\frac{|\nabla f(z)|}{|x-z|^{d-1}}
 \,\mathrm d\lambda_d(z).
\end{align*}
Druga nierówność wynika z obliczenia, dla $s=|x-z|>0$,
\[
 \int_{s/(2R)}^1 t^{-d-1}\,\mathrm d\lambda_1(t)
 =\frac{(2R/s)^d-1}{d}\le\frac{(2R)^d}{ds^d}.
\]
Ponieważ $|B|=\omega_dR^d$, gdzie
$\omega_d=\lambda_d(B_1(0))$, można przyjąć
$C_d=2^d/(d\omega_d)$. Punkt $z=x$ pomijamy jako zbiór
miary zero. Jądro $|x-z|^{1-d}$ jest lokalnie całkowalne:
całkowanie w biegunowych na kuli o promieniu $2R$
daje $d\omega_d\int_0^{2R}1\,\mathrm d\lambda_1(s)$.
W szczególności prawa strona jest skończona
dla rozważanej funkcji klasy $C^1$.
\end{proof}

\begin{twierdzenie}[Morrey na $\R^d$]
\label{thm:pde10-Morrey}
Niech $d\ge2$ i $d<p<\infty$. Każda klasa
$u\in W^{1,p}(\R^d)$ ma jednoznacznego ciągłego
reprezentanta $\widetilde u$ i istnieje stała $C(d,p)$ taka, że
\[
 |\widetilde u(x)-\widetilde u(y)|
 \le C(d,p)|x-y|^{1-d/p}\|\nabla u\|_{L^p(\R^d)}
 \quad(x,y\in\R^d).
\]
Dodatkowo $\|\widetilde u\|_\infty\le
 C(d,p)\|u\|_{W^{1,p}(\R^d)}$.
Dla $p=\infty$ każda klasa $u\in W^{1,\infty}(\R^d)$
ma jedyny ciągły reprezentant globalnie lipschitzowski
z $\operatorname{Lip}\widetilde u\le\|\nabla u\|_\infty$
oraz $\|\widetilde u\|_\infty=\|u\|_{L^\infty}$.
\end{twierdzenie}
\begin{proof}
Najpierw rozważmy $u\in C_c^\infty(\R^d)$. Ustalmy
$x\ne y$ i oznaczmy $r=|x-y|$.
Wybieramy kulę $B=B_{2r}(x)$, która zawiera oba punkty.
Z lematu~\ref{lem:pde10-point-average} otrzymujemy
$|u(x)-u(y)|\le|u(x)-u_B|+|u(y)-u_B|$.
Każdy składnik szacujemy nierównością Höldera
z wykładnikami $p$ oraz $p'=p/(p-1)$.
Dla $z_0=x$ lub $z_0=y$ liczymy
\[
 \left(\int_B |z-z_0|^{-(d-1)p'}
 \,\mathrm d\lambda_d(z)\right)^{1/p'}
 \le C(d,p)r^{1-d/p}.
\]
Istotnie, $B\subset B_{4r}(z_0)$, a w biegunowych
całka jest ograniczona przez
$C_d\int_0^{4r}s^{d-1-(d-1)p'}\,\mathrm d\lambda_1(s)$.
Mamy
\[
 d-(d-1)p'=\frac{p-d}{p-1}>0,\qquad
 \frac{d-(d-1)p'}{p'}=1-\frac dp.
\]
Całka potęgowa jest więc skończona i po wyciągnięciu
$p'$-tego pierwiastka daje dokładnie czynnik
$r^{1-d/p}$. Z lematu i Höldera otrzymujemy
\[
 |u(z_0)-u_B|\le C(d,p)r^{1-d/p}\|\nabla u\|_{L^p(B)}.
\]
Sumując dla $z_0=x,y$, dostajemy oszacowanie z tezy;
dla $x=y$ obie strony są zerowe.

Aby kontrolować wartość w jednym punkcie,
uśredniamy $u$ na kuli jednostkowej $B_1(x)$.
Stosujemy bezpośrednio lemat~\ref{lem:pde10-point-average}
na $B_1(x)$ i Höldera. Rachunek całki jądra dla promienia
$1$ daje $|u(x)-u_{B_1(x)}|\le C(d,p)\|\nabla u\|_p$,
a z Höldera $|u_{B_1(x)}|\le\omega_d^{-1/p}\|u\|_p$.
Dodając oszacowania i biorąc supremum po $x$, uzyskujemy
$\|u\|_\infty\le C(d,p)\|u\|_{W^{1,p}}$,
z uwzględnieniem równoważności norm gradientu.

Dla ogólnego $u\in W^{1,p}(\R^d)$ wybieramy
gładkie funkcje o zwartym nośniku $u_n\to u$ w $W^{1,p}$
na mocy gęstości z sekcji~\ref{sec:pde9-density}.
Oszacowanie supremum zastosowane do $u_n-u_m$ daje
ciąg Cauchy'ego w normie jednostajnej. Jego granica
$\widetilde u$ jest ograniczona i ciągła. Wybieramy
podciąg z $\|u_{n_j}-u\|_p^p<2^{-j}$. Z Tonellego
$\sum_j|u_{n_j}-u|^p<\infty$ prawie wszędzie, więc
$u_{n_j}\to u$ prawie wszędzie. Jednostajna granica
$\widetilde u$ reprezentuje zatem klasę $u$.
Dwa ciągłe reprezentanty równe prawie wszędzie są równe
wszędzie: niezerowa różnica w punkcie byłaby niezerowa
na otwartej kuli dodatniej miary.
W oszacowaniu przyrostów dla $u_n$ przechodzimy
do granicy dla każdego $x,y$, ponieważ zbieżność w
$W^{1,p}$ daje zbieżność norm gradientów. Analogicznie
przechodzimy do granicy w oszacowaniu supremum.
Jednoznaczność reprezentanta zapewnia liniowość
przyporządkowania $u\mapsto\widetilde u$.

Dla $p=\infty$ niech $u_\varepsilon=\rho_\varepsilon*u$
na $\R^d$, z nieujemnym jądrem o całce $1$ z rozdziału
~\ref{chap:pde-sobolev}. Z twierdzenia
~\ref{thm:pde9-mollify}, stosowanego lokalnie przy $p=1$,
$\nabla u_\varepsilon=\rho_\varepsilon*(\nabla u)$.
Nieujemność jądra i nierówność trójkąta dla normy
euklidesowej dają w każdym punkcie
\[
 |u_\varepsilon(x)|\le\|u\|_\infty,\qquad
 |\nabla u_\varepsilon(x)|\le\|\nabla u\|_\infty.
\]
Całkowanie po odcinku daje
$|u_\varepsilon(x)-u_\varepsilon(y)|
 \le\|\nabla u\|_\infty|x-y|$.
Z lokalnej zbieżności w $L^1$ wybieramy
$\varepsilon_n\downarrow0$ tak, aby
$\|u_{\varepsilon_n}-u\|_{L^1(B_n(0))}<2^{-n}$.
Dla ustalonego $m$ suma tych błędów na $B_m(0)$
od $n=m$ jest skończona; Tonelli daje zbieżność
punktową prawie wszędzie na $B_m(0)$.
Po przeliczalnej sumie zbiorów wyjątkowych otrzymujemy
zbiór $S\subset\R^d$ pełnej miary, na którym
$u_{\varepsilon_n}\to u$. Zbiór $S$ jest gęsty,
gdyż każda niepusta kula ma dodatnią miarę.
Przejście do granicy dla $x,y\in S$ daje na $S$
to samo oszacowanie lipschitzowskie i ograniczenie
modułu przez $\|u\|_\infty$. Dla każdego $x\in\R^d$
wybieramy $x_j\in S$, $x_j\to x$, i definiujemy
$\widetilde u(x)=\lim_j u(x_j)$.
Oszacowanie przyrostów daje istnienie granicy
i jej niezależność od wyboru ciągu, a po przejściu
do granicy również globalną stałą Lipschitza.
Na $S$ reprezentant zgadza się z $u$.
Jego supremum jest równe istotnemu supremum dzięki
ciągłości, a jednoznaczność wynika jak wyżej.
Nie korzystaliśmy ze zbieżności w $W^{1,\infty}$.
\end{proof}

\begin{wniosek}[Morrey na obszarze i wyższe rzędy]
\label{cor:pde10-Morrey-domain}
Jeśli $\Omega\subset\R^d$ jest obszarem ograniczonym Lipschitzowskim i $d<p\le\infty$, to
$W^{1,p}(\Omega)\hookrightarrow C^{0,1-d/p}(\overline\Omega)$,
z naturalną interpretacją wykładnika $1$ dla $p=\infty$.
Jeśli $k\ge2$ i $u\in W^{k,p}(\Omega)$, $d<p\le\infty$,
to pochodne słabe $\partial^\alpha u$ dla $|\alpha|\le k-1$
mają ciągłe reprezentanty H\"olderowskie wykładnika
$1-d/p$; stąd $u$ ma reprezentanta w $C^{k-1,1-d/p}$.
\end{wniosek}
\begin{proof}
Dla $k=1$, $d\ge2$ bierzemy reprezentanta ciągłego
$Eu$ z twierdzenia~\ref{thm:pde10-Morrey}, gdzie $E$
jest operatorem z lematu~\ref{lem:pde10-extension}.
Jego ograniczenie do $\overline\Omega$ reprezentuje $u$
na $\Omega$. Norma supremum i półnorma Höldera
na $\overline\Omega$ nie przekraczają odpowiednich norm
na $\R^d$. Otrzymujemy, przy $\alpha=1-d/p$,
\[
 \|\widetilde u\|_{C^{0,\alpha}(\overline\Omega)}
 \le C(d,p)\|Eu\|_{W^{1,p}(\R^d)}
 \le C(\Omega,p)\|u\|_{W^{1,p}(\Omega)}.
\]
Dla $p=\infty$ używamy $\alpha=1$ i oszacowania
Lipschitza wraz z normą supremum. Dla $d=1$ obszar
jest ograniczonym przedziałem; stosujemy twierdzenie
~\ref{thm:pde9-ac} i wniosek~\ref{cor:pde9-trace-oneD}
oraz oszacowanie przyrostów z tego twierdzenia i Höldera
(dla $p=\infty$ z normą supremum pochodnej).
Jednoznaczność wynika z równości prawie wszędzie
funkcji ciągłych na $\Omega$ i z gęstości $\Omega$
w $\overline\Omega$. Liniowość wynika z jednoznaczności.

Dla $k\ge2$ każdy $\partial^\beta u$, $|\beta|\le k-1$,
należy do $W^{1,p}(\Omega)$, ponieważ jego słabe
pochodne pierwszego rzędu są pochodnymi $u$ rzędu
$|\beta|+1\le k$. Niech $v_\beta$ będzie jego ciągłym
reprezentantem na $\overline\Omega$. Wszystkie te
funkcje spełniają oszacowania Höldera i supremum
przez stałą razy $\|u\|_{W^{k,p}(\Omega)}$.
Dla $|\beta|\le k-2$ i $x\in\Omega$ wybieramy kulę
$B_{2r}(x)\Subset\Omega$. Na $B_r(x)$ lokalne sploty
$u_\varepsilon$ oraz ich pochodne do rzędu $k-1$
są splotami odpowiednich $v_\gamma$ z jądrem, na mocy
twierdzenia~\ref{thm:pde9-mollify} (dla $p=\infty$
stosujemy lokalnie przypadek $p=1$).
Z jednostajnej ciągłości $v_\gamma$ na
$\overline{B_{2r}(x)}$ otrzymujemy dla $\varepsilon<r$
\[
 \sup_{z\in B_r(x)}
 |\partial^\gamma u_\varepsilon(z)-v_\gamma(z)|
 \le\sup_{\substack{z\in B_r(x)\\|h|\le\varepsilon}}
 |v_\gamma(z-h)-v_\gamma(z)|\longrightarrow0.
\]
Dla $|t|<r$ podstawowe twierdzenie rachunku całkowego
daje
\[
 \partial^\beta u_\varepsilon(x+te_i)
 -\partial^\beta u_\varepsilon(x)
 =\int_0^t\partial^{\beta+e_i}u_\varepsilon(x+se_i)
             \,\mathrm d\lambda_1(s).
\]
Przechodząc do granicy jednostajnej po obu stronach,
otrzymujemy tę samą tożsamość dla $v_\beta$ oraz
$v_{\beta+e_i}$. Dzieląc przez $t\ne0$ i korzystając
z ciągłości $v_{\beta+e_i}$ przy $t\to0$, mamy
klasyczną pochodną $\partial_i v_\beta=v_{\beta+e_i}$.
Wszystkie te pochodne są ciągłe, więc indukcyjnie
$v_0\in C^{k-1}(\Omega)$, a jego klasyczne pochodne
przedłużają się ciągle i hölderowsko na
$\overline\Omega$. To właśnie oznacza tu
$C^{k-1,\alpha}(\overline\Omega)$; nie wymaga
istnienia klasycznego przedłużenia wyższego rzędu
poza obszar.
\end{proof}

\begin{uwaga}[Wymiar $d=1$]
Na ograniczonym przedziale wyniki są silniejsze:
$W^{1,1}(a,b)$ ma reprezentanta absolutnie ciągłego,
a dla $1<p<\infty$ nierówność Höldera i twierdzenie
~\ref{thm:pde9-ac} dają
$|u(x)-u(y)|\le |x-y|^{1-1/p}\|u'\|_{L^p(a,b)}$.
Dla $p=\infty$ otrzymujemy oszacowanie lipschitzowskie
z wykładnikiem $1$. Dla $p=1$ nie ma wspólnego
dodatniego wykładnika Höldera (przykłady $u(x)=x^\theta$
na $(0,1)$, $0<\theta<1$, omówione w rozdziale
~\ref{chap:pde-sobolev}).
Nie wolno utożsamiać tego przypadku z uogólnieniem wzoru
$p^*=dp/(d-p)$ na $d=1$.
\end{uwaga}

\section{Zwartość: od średnich po kostkach do twierdzenia Rellicha}
\label{sec:pde10-Rellich}

\begin{lemat}[Aproksymacja funkcjami stałymi na kostkach]
\label{lem:pde10-cube-approx}
Niech $Q\subset\R^d$ będzie kostką o krawędzi $L$, a
$1\le p<\infty$. Dla każdego $N\in\N$ podzielmy $Q$
na $N^d$ równych kostek $Q_j$ o krawędzi $h=L/N$,
a dla $u\in L^1(Q)$ połóżmy
$P_Nu=\sum_j u_{Q_j}\chi_{Q_j}$, gdzie
$u_{Q_j}=|Q_j|^{-1}\int_{Q_j}u\,\mathrm d\lambda_d$.
Dla wszystkich $u\in W^{1,p}(Q)$ zachodzi
\[
 \|u-P_Nu\|_{L^p(Q)}\le dh\|\nabla u\|_{L^p(Q)}.
\]
\end{lemat}
\begin{proof}
Stosujemy nierówność Poincar\'ego--Wirtingera
z twierdzenia~\ref{thm:pde9-cube-Poincare} osobno na
każdej $Q_j$:
$\|u-u_{Q_j}\|_{L^p(Q_j)}\le
 h\sum_{i=1}^d\|\partial_i u\|_{L^p(Q_j)}$.
Ponieważ $(\sum_i a_i)^p\le d^{p-1}\sum_i a_i^p$
dla $a_i\ge0$, po podniesieniu do potęgi $p$
i zsumowaniu po rozłącznych kostkach otrzymujemy
$\|u-P_Nu\|_p^p\le d^{p-1}h^p\sum_i\|\partial_i u\|_p^p$.
Po wyciągnięciu pierwiastka mamy w istocie stałą
$d^{1-1/p}$; oszacowanie przez $d$ jest wystarczające.
\end{proof}

\begin{twierdzenie}[Rellich--Kondrachov: podstawowa zwartość]
\label{thm:pde10-Rellich-p}
Niech $\Omega\subset\R^d$ będzie obszarem ograniczonym Lipschitzowskim i $1\le p<\infty$. Wtedy zanurzenie
$W^{1,p}(\Omega)\hookrightarrow\hookrightarrow L^p(\Omega)$
jest zwarte.
\end{twierdzenie}
\begin{proof}
Niech $(u_n)$ będzie ograniczony w $W^{1,p}(\Omega)$.
Przedłużamy każdą funkcję operatorem
$E$ z lematu~\ref{lem:pde10-extension}. Wszystkie
$Eu_n$ mają nośniki zawarte w jednej ustalonej kostce $Q$
i są ograniczone w $W^{1,p}(Q)$ przez stałą $M$.

Ustalmy $\varepsilon>0$. Z lematu~\ref{lem:pde10-cube-approx}
wybieramy $N$ tak duże, aby
$\|Eu_n-P_NEu_n\|_{L^p(Q)}<\varepsilon/3$
jednocześnie dla wszystkich $n$.
Funkcje $P_NEu_n$ należą do skończeniewymiarowej przestrzeni
funkcji stałych na kostkach $Q_j$, a ich normy są
jednostajnie ograniczone przez $M$ (Jensen daje
$\|P_Nf\|_p\le\|f\|_p$). Zbiór ograniczony w przestrzeni
skończeniewymiarowej jest całkowicie ograniczony,
więc ma skończoną $\varepsilon/3$-sieć.
Z nierówności trójkąta wynika, że zbiór $\{Eu_n\}$
ma skończoną $\varepsilon$-sieć w $L^p(Q)$.

Dla pełnej ścisłości, wybierając kolejno sieci
dla $\varepsilon=2^{-m}$ i zagnieżdżone podciągi,
konstrukcja przekątniowa daje podciąg $(Eu_{n_k})$
Cauchy'ego w $L^p(Q)$; zupełność $L^p$
zapewnia jego zbieżność. Ograniczenie do $\Omega$
nie zwiększa normy, więc $u_{n_k}$ zbiega w
$L^p(\Omega)$. Nie wykorzystaliśmy
refleksywności $W^{1,p}$, dlatego dowód działa również
przy $p=1$.
\end{proof}

\begin{lemat}[Interpolacja elementarna norm $L^p$]
\label{lem:pde10-interpolation}
Niech $1\le a<q<b<\infty$, a $f\in L^a(\Omega)\cap L^b(\Omega)$.
Jeśli $0<\theta<1$ spełnia $1/q=\theta/a+(1-\theta)/b$,
to
$\|f\|_{L^q}\le\|f\|_{L^a}^\theta
\|f\|_{L^b}^{1-\theta}$.
\end{lemat}
\begin{proof}
Piszemy $|f|^q=(|f|^a)^{q\theta/a}
(|f|^b)^{q(1-\theta)/b}$.
Dwa wykładniki sumują się do $1$.
Oznaczmy $J_r=\int_\Omega|f(x)|^r\,\mathrm d\lambda_d(x)$
dla $r\in\{a,q,b\}$. Nierówność H\"oldera z wykładnikami
$a/(q\theta)$ i $b/(q(1-\theta))$ daje
\[
 J_q\le J_a^{q\theta/a}J_b^{q(1-\theta)/b},
\]
co kończy dowód.
\end{proof}

\begin{twierdzenie}[Rellich--Kondrachov: wykładniki podkrytyczne]
\label{thm:pde10-Rellich-q}
Niech $\Omega\subset\R^d$ będzie obszarem ograniczonym Lipschitzowskim, $d\ge2$.
\begin{enumerate}[label=\textnormal{(\roman*)}]
\item Jeśli $1\le p<d$, to
$W^{1,p}(\Omega)\hookrightarrow\hookrightarrow L^q(\Omega)$
dla $1\le q<p^*=dp/(d-p)$.
\item Jeśli $p=d$, to to samo zanurzenie jest zwarte dla
każdego skończonego $q\ge1$.
\item Jeśli $d<p\le\infty$, to
$W^{1,p}(\Omega)\hookrightarrow\hookrightarrow
C^{0,\beta}(\overline\Omega)$
dla każdego $0\le\beta<1-d/p$; przy $p=\infty$ wolno
przyjąć dowolne $\beta<1$.
\end{enumerate}
\end{twierdzenie}
\begin{proof}
W (i) niech $(u_n)$ będzie ograniczony w $W^{1,p}$.
Z twierdzenia~\ref{thm:pde10-Rellich-p} wybieramy
podciąg Cauchy'ego w $L^p$. Twierdzenie
\ref{thm:pde10-domain-embed} daje jednostajne
ograniczenie norm $L^{p^*}$, zatem różnice
$u_n-u_m$ mają jednostajnie ograniczone normy $L^{p^*}$.
Dla $p<q<p^*$ lemat~\ref{lem:pde10-interpolation}
implikuje, że podciąg jest Cauchy'ego w $L^q$.
Dla $q\le p$ wystarcza nierówność H\"oldera na
zbiorze o skończonej mierze. Tak uzyskujemy (i).

W (ii) jeśli $q>p=d$, wybieramy ustalone skończone
$r>q$. Zanurzenie $W^{1,d}\hookrightarrow L^r$ daje
ograniczenie norm $L^r$, a interpolacja między $L^d$
i $L^r$ daje zbieżność w $L^q$. Dla $q\le d$
postępujemy jak wcześniej.

W (iii) nierówność Morreya daje jednostajne ograniczenie
norm $C^{0,\alpha}(\overline\Omega)$, gdzie
$\alpha=1-d/p>0$ (albo $\alpha=1$ dla $p=\infty$).
Rodzina jest więc jednostajnie ograniczona i jednakowo
ciągła na zwartym $\overline\Omega$.
Twierdzenie Arzel\`a--Ascoliego zapewnia podciąg zbieżny
jednostajnie do ciągłej funkcji $u$.
Jeśli $f_n=u_n-u$, to $\|f_n\|_\infty\to0$
i $[f_n]_{C^{0,\alpha}}\le C$.
Dla $0<\beta<\alpha$ i każdego $\delta>0$ mamy
\[
 [f_n]_{C^{0,\beta}}
 \le\max\{2\|f_n\|_\infty\delta^{-\beta},
 C\delta^{\alpha-\beta}\}.
\]
Najpierw wybieramy $\delta$ małe, by drugi wyraz
był mały, potem $n$ duże, by pierwszy też był mały.
To dowodzi zbieżności w $C^{0,\beta}$.
Dla $\beta=0$ oznacza ona zbieżność jednostajną.
\end{proof}

\begin{kontrprzyklad}[Brak zwartości dla wykładnika krytycznego]
\label{ex:pde10-noncompact-critical}
Niech $1\le p<d$, $x_0\in\Omega$ oraz
$\varphi\in C_c^\infty(B_1(0))$ będzie niezerowa.
Funkcje $u_\varepsilon(x)=\varepsilon^{1-d/p}
\varphi((x-x_0)/\varepsilon)$, dla dostatecznie małych
$\varepsilon>0$, mają jednostajnie ograniczone normy
$W^{1,p}(\Omega)$ i stałą dodatnią normę $L^{p^*}(\Omega)$.
Jednocześnie $u_\varepsilon\to0$ prawie wszędzie, gdy
$\varepsilon\downarrow0$. Gdyby istniał podciąg zbieżny
silnie w $L^{p^*}$, jego granica musiałaby być równa zeru,
co przeczyłoby stałości norm. Dlatego zanurzenie przy
$q=p^*$ nie jest zwarte.
\end{kontrprzyklad}

\begin{kontrprzyklad}[Brak zwartości na całej $\R^d$]
\label{ex:pde10-translation}
Niech $\varphi\in C_c^\infty(B_1(0))$ będzie niezerowa
i $u_n(x)=\varphi(x-3ne_1)$.
Wówczas $\|u_n\|_{W^{1,p}(\R^d)}$ jest stała dla
każdego $1\le p<\infty$, lecz dla $m\ne n$ nośniki są
rozłączne, więc
$\|u_n-u_m\|_p^p=2\|\varphi\|_p^p$.
Żaden podciąg nie może być Cauchy'ego w $L^p$.
Ograniczoność obszaru w twierdzeniu Rellicha nie jest
zatem przypadkowym warunkiem technicznym.
\end{kontrprzyklad}

\section[Przypadek graniczny w zanurzeniach Sobolewa]{Przypadek graniczny $p=d$: dlaczego nie dostajemy $L^\infty$?}
\label{sec:pde10-Moser}

\begin{kontrprzyklad}[Ciąg koncentracyjny typu Mosera]
\label{ex:pde10-Moser}
Niech $d\ge2$, $B=B_1(0)$ oraz $0<\varepsilon<e^{-1}$.
Połóżmy $L=\log(1/\varepsilon)$ i
\[
 u_\varepsilon(x)=
 \begin{cases}
 L^{1-1/d}, & |x|\le\varepsilon,\\
 L^{-1/d}\log(1/|x|), & \varepsilon<|x|<1,\\
 0,&|x|\ge1.
 \end{cases}
\]
Funkcja jest ciągła i należy do $W^{1,d}(\R^d)$.
W istocie dla $\varepsilon<|x|<1$ mamy
$|\nabla u_\varepsilon(x)|=L^{-1/d}/|x|$, skąd
\[
 \|\nabla u_\varepsilon\|_{L^d(\R^d)}^d
 =|S^{d-1}|L^{-1}\int_\varepsilon^1\frac{\mathrm dr}{r}
 =|S^{d-1}|.
\]
Natomiast $\|u_\varepsilon\|_{L^d}$ jest jednostajnie
ograniczona: na $\varepsilon<|x|<1$ wystarcza
$L^{-1}\int_0^\infty t^de^{-dt}\,\mathrm dt<\infty$,
a na kuli $B_\varepsilon$ wkład wynosi
$\omega_d\varepsilon^dL^{d-1}$ i pozostaje ograniczony.
Pomimo tego
$\|u_\varepsilon\|_\infty=L^{1-1/d}\to\infty$.
Nie istnieje więc ciągłe zanurzenie
$W^{1,d}(B)\hookrightarrow L^\infty(B)$.
\end{kontrprzyklad}

\begin{uwaga}[Ostrość wykładnika Morreya]
Niech $p>d$ i $\alpha=1-d/p$. Dla tej samej skalowanej
rodziny $u_\varepsilon(x)=\varepsilon^\alpha
\varphi((x-x_0)/\varepsilon)$ normy $W^{1,p}$ są
ograniczone, lecz dla $\beta>\alpha$ ilorazy H\"oldera
rosną jak $\varepsilon^{\alpha-\beta}$.
Nie można więc ogólnie zastąpić wykładnika Morreya
większym wykładnikiem, nawet na kuli.
\end{uwaga}

\section{Ślady na brzegu Lipschitzowskim}
\label{sec:pde10-traces}

\begin{definicja}[Miara i przestrzeń $L^p$ na brzegu]
Jeśli $\Omega\subset\R^d$ jest obszarem ograniczonym Lipschitzowskim, przyjmujemy na $\partial\Omega$ miarę
$\mathcal H^{d-1}|_{\partial\Omega}$, czyli ograniczenie
$(d-1)$-wymiarowej miary Hausdorffa. Na każdym lokalnym
wykresie $x_d=\phi(x')$ miara ta ma element
$\sqrt{1+|\nabla\phi(x')|^2}\,\mathrm d\lambda_{d-1}(x')$
prawie wszędzie. Przestrzeń $L^p(\partial\Omega)$
rozumiemy względem tej miary.
\end{definicja}

\begin{twierdzenie}[Twierdzenie o śladzie]
\label{thm:pde10-trace}
Niech $\Omega\subset\R^d$ będzie obszarem ograniczonym Lipschitzowskim, $1\le p<\infty$. Istnieje dokładnie jeden
ograniczony operator liniowy
$\operatorname{Tr}:W^{1,p}(\Omega)\to L^p(\partial\Omega)$,
który dla każdej $u\in C^\infty(\overline\Omega)$
(gładkiej w otoczeniu $\overline\Omega$) daje zwykłą
restrykcję $u|_{\partial\Omega}$. Ponadto
\[
 \|\operatorname{Tr}u\|_{L^p(\partial\Omega)}
 \le C(\Omega,p)\|u\|_{W^{1,p}(\Omega)}.
\]
\end{twierdzenie}
\begin{proof}
Najpierw udowodnimy oszacowanie na płaskim brzegu.
Dla $w\in C^1(\R^{d-1}\times[0,a])$ o odpowiedniej
całkowalności, $t\in(0,a)$, podstawowe twierdzenie
rachunku daje
$|w(x',0)|\le|w(x',t)|+
\int_0^a|\partial_tw(x',s)|\,\mathrm d\lambda_1(s)$.
Całkujemy nierówność po $t\in(0,a)$, dzielimy przez $a$
i stosujemy H\"oldera. Otrzymujemy punktowo
\[
 |w(x',0)|\le a^{-1/p}\|w(x',\cdot)\|_{L^p(0,a)}
 +a^{1-1/p}\|\partial_tw(x',\cdot)\|_{L^p(0,a)}.
\]
Po podniesieniu do potęgi $p$ i scałkowaniu po $x'$
dostajemy
\[
 \|w(\cdot,0)\|_{L^p(\R^{d-1})}^p
 \le 2^{p-1}\bigl(a^{-1}\|w\|_{L^p}^p
 +a^{p-1}\|\partial_tw\|_{L^p}^p\bigr).
\]
W lokalnym wykresie brzegu prostujemy go odwzorowaniem
$\Phi(x',t)=(x',t+\phi(x'))$. Funkcję $\chi u$,
gdzie $\chi$ jest odcięciem wewnątrz mniejszego cylindra,
przenosimy do płaskiej warstwy. Zmiana zmiennych
ma jakobian $1$, a $\partial_t[(\chi u)\circ\Phi]$
zgadza się z $\partial_{x_d}(\chi u)\circ\Phi$.
Na wykresie element powierzchniowy jest ograniczony
z góry przez $\sqrt{1+\operatorname{Lip}(\phi)^2}$
razy $\mathrm d\lambda_{d-1}$.
Otrzymujemy więc lokalne oszacowanie śladu przez
$\|\chi u\|_{W^{1,p}(\Omega)}$.
Skończone pokrycie brzegu i podział jedności pozwalają
zsumować te nierówności; stąd globalne oszacowanie
na gładkich $u$.

Aby zdefiniować ślad ogólnej funkcji, wybieramy
ciąg restrykcji funkcji $C_c^\infty(\R^d)$ zbieżny
w $W^{1,p}(\Omega)$, którego istnienie wykazano
we wniosku~\ref{cor:pde10-density-boundary}.
Otrzymane oszacowanie pokazuje, że restrykcje przybliżeń
tworzą ciąg Cauchy'ego w $L^p(\partial\Omega)$;
jego granica nie zależy od wyboru przybliżeń.
Przedłużenie operatora jest liniowe i ograniczone,
a gęstość funkcji gładkich aż do brzegu zapewnia
jednoznaczność.
\end{proof}

\begin{definicja}[Przestrzeń danych brzegowych $H^{1/2}(\partial\Omega)$]
\label{def:pde10-halftrace}
Niech $\Omega\subset\R^d$ będzie ograniczonym obszarem Lipschitzowskim,
a $\operatorname{Tr}:H^1(\Omega)\to L^2(\partial\Omega)$
operatorem z twierdzenia~\ref{thm:pde10-trace}. Definiujemy
\[
 H^{1/2}(\partial\Omega):=\operatorname{Tr}(H^1(\Omega)),\qquad
 \|g\|_{H^{1/2}(\partial\Omega)}
 :=\inf_{\substack{G\in H^1(\Omega)\\\operatorname{Tr}G=g}}
 \|G\|_{H^1(\Omega)}.
\]
Jest to definicja przez przestrzeń śladów, z normą ilorazową.
Na ograniczonym brzegu Lipschitzowskim jest ona równoważna
zwyczajowej definicji przez ułamkowe seminormy Gagliarda;
dowód tej identyfikacji wymaga oddzielnego twierdzenia o śladach.
\end{definicja}

\begin{lemat}[Liniowy operator podnoszenia śladu]
\label{lem:pde10-halftrace-rightinverse}
Dla każdego ograniczonego obszaru Lipschitzowskiego $\Omega$
istnieje liniowy, ograniczony operator
$E_{\operatorname{Tr}}:H^{1/2}(\partial\Omega)\to H^1(\Omega)$
taki, że $\operatorname{Tr}E_{\operatorname{Tr}}g=g$ oraz
$\|E_{\operatorname{Tr}}g\|_{H^1(\Omega)}
=\|g\|_{H^{1/2}(\partial\Omega)}$.
\end{lemat}
\begin{proof}
Niech $H=H^1(\Omega)$ i $N=\ker\operatorname{Tr}$.
Ponieważ $\operatorname{Tr}:H\to L^2(\partial\Omega)$
jest ciągły, $N$ jest domkniętą podprzestrzenią Hilberta.
Mamy więc rozkład ortogonalny $H=N\oplus N^\perp$.
Dla $g\in\operatorname{Tr}(H)$ wybieramy dowolny $G\in H$
o śladzie $g$ i definiujemy
$E_{\operatorname{Tr}}g=P_{N^\perp}G$.
Jeżeli $G_1$ i $G_2$ mają ten sam ślad, to
$G_1-G_2\in N$, a zatem ich rzuty na $N^\perp$ są równe.
Konstrukcja jest więc niezależna od wyboru $G$ i liniowa.
Ponadto $G-P_{N^\perp}G\in N$, skąd
$\operatorname{Tr}E_{\operatorname{Tr}}g=g$.
Każde inne podniesienie jest postaci
$E_{\operatorname{Tr}}g+n$ dla $n\in N$, więc
z twierdzenia Pitagorasa
$\|E_{\operatorname{Tr}}g+n\|_H^2
=\|E_{\operatorname{Tr}}g\|_H^2+\|n\|_H^2$.
To dowodzi równości norm i ograniczoności operatora.
W szczególności norma ilorazowa czyni
$H^{1/2}(\partial\Omega)$ przestrzenią Hilberta.
\end{proof}

\begin{twierdzenie}[Zerowy ślad i domknięcie funkcji testowych]
\label{thm:pde10-trace-zero}
Niech $\Omega\subset\R^d$ będzie obszarem ograniczonym Lipschitzowskim, $1\le p<\infty$. Wówczas
\[
 W_0^{1,p}(\Omega)=\ker\operatorname{Tr}.
\]
\end{twierdzenie}
\begin{proof}
Jeśli $u\in W_0^{1,p}(\Omega)$, to z definicji jest
granicą w $W^{1,p}$ funkcji $C_c^\infty(\Omega)$.
Te funkcje mają zerowy ślad, więc przez ciągłość
$\operatorname{Tr}u=0$.

Dla implikacji odwrotnej niech $\operatorname{Tr}u=0$.
Rozważmy lokalny fragment brzegu wyprostowany do
półwarstwy $t>0$. Z twierdzenia o absolutnej ciągłości
na prawie wszystkich pionowych odcinkach dla funkcji
$W^{1,p}$ wynika, że dla prawie wszystkich $x'$
wartość brzegowa $w(x',0)$ istnieje i pokrywa się
ze śladem (dowodzi się tego przez gęstość w normie
$W^{1,p}$ i jednowymiarowe oszacowanie z poprzedniego
dowodu). Jest zatem równa zero. Przedłużenie $w$
przez zero na $t<0$ ma słabą pochodną normalną
równą przedłużeniu $\partial_tw$ przez zero:
podczas całkowania przez części w $t$ składnik brzegowy
$w(x',0)$ znika. Pochodne styczne również
przedłużają się przez zero, gdyż granica całkowania
w tych współrzędnych jest niezależna od $x'$.

Wybierzmy funkcję odcinającą $\chi$ o nośniku w
mniejszym cylindrze oraz dodatni kierunek pionowy
skierowany do wnętrza obszaru. Przesuwamy
przedłużenie $\chi w$ o odległość $a\varepsilon$
do wnętrza, a następnie wygładzamy jądrem o promieniu
$\varepsilon$; wybieramy stałą $a>1$ tak, aby wygładzenie
miało nośnik po stronie $t>0$. Translacja jest ciągła
w $W^{1,p}(\R^d)$ dla $p<\infty$, a splot
z przybliżoną jedynką zbiega w tej samej normie.
Wszystkie uzyskane funkcje są zatem gładkie, wspierają
się we wnętrzu odpowiedniego fragmentu obszaru
i zbiegają do $\chi w$ w $W^{1,p}$.

Po powrocie do pierwotnych współrzędnych na wykresie
Lipschitzowskim gładkość nie musi być zachowana.
Dlatego konstrukcję wykonujemy geometrycznie przed
prostowaniem: przesuwamy fragment $\chi u$ w stałym
kierunku pionowym oryginalnego cylindra na odległość
$a\varepsilon$, gdzie $a>1+\operatorname{Lip}\phi$,
a następnie wygładzamy w oryginalnych współrzędnych
promieniem $\varepsilon$. Warunek na $a$ gwarantuje,
że całe kule używane w splocie pozostają po stronie
wewnętrznej grafu. Zbieżność w $W^{1,p}$ wynika
z własności translacji w $\R^d$ i ograniczoności
operatorów mnożenia przez odcięcia.
Części podziału jedności wspierające się ściśle
wewnątrz $\Omega$ wygładzamy bez przesunięć.
Sumując skończenie wiele aproksymacji, otrzymujemy
ciąg w $C_c^\infty(\Omega)$ zbieżny do $u$
w $W^{1,p}(\Omega)$. To kończy dowód.
\end{proof}

\begin{wniosek}[Nierówność Poincar\'ego na obszarze]
\label{cor:pde10-Poincare}
Niech $\Omega\subset\R^d$ będzie obszarem ograniczonym Lipschitzowskim, $1\le p<\infty$.
Istnieje stała $C=C(\Omega,p)$ taka, że
\[
 \|u-u_\Omega\|_{L^p(\Omega)}
 \le C\|\nabla u\|_{L^p(\Omega)},\qquad
 u_\Omega=|\Omega|^{-1}\int_\Omega u\,\mathrm d\lambda_d.
\]
Dla $u\in W_0^{1,p}(\Omega)$ istnieje również stała
$C'$ z $\|u\|_p\le C'\|\nabla u\|_p$.
\end{wniosek}
\begin{proof}
Gdyby pierwsza nierówność była fałszywa, istniałyby
$u_n\in W^{1,p}(\Omega)$ takie, że
$(u_n)_\Omega=0$, $\|u_n\|_p=1$ i
$\|\nabla u_n\|_p\to0$ (odejmujemy średnią i
normalizujemy przez normę $L^p$).
Ciąg jest ograniczony w $W^{1,p}$, więc twierdzenie
\ref{thm:pde10-Rellich-p} daje podciąg zbieżny w
$L^p$ do pewnego $u$ o normie $1$ i średniej $0$.
Z zamkniętości pochodnej dystrybucyjnej
(wniosek~\ref{cor:pde9-closed}) otrzymujemy
$\nabla u=0$. Na obszarze funkcja lokalnie
całkowalna o zerowych pochodnych dystrybucyjnych
jest stała prawie wszędzie; jej średnia jest zerowa,
więc $u=0$, sprzeczność.
Dla drugiej nierówności przedłużamy funkcję z
$W_0^{1,p}(\Omega)$ przez zero na kostkę $Q\supset\Omega$
i stosujemy wniosek~\ref{cor:pde9-Poincare-zero}.
\end{proof}

% Audyt rozdziału 10: jawny krok techniczny używany w dalszej części tomu.
\begin{wniosek}[Przedłużenie przez zero funkcji o zerowym śladzie]
\label{cor:pde10-zeroextension}
Niech $\Omega\subset\R^d$ będzie obszarem ograniczonym Lipschitzowskim, $1\le p<\infty$ i
$u\in W^{1,p}(\Omega)$ z $\operatorname{Tr}u=0$.
Niech $E_0u$ oznacza funkcję równą $u$ prawie wszędzie
na $\Omega$ i $0$ poza $\Omega$. Wówczas
$E_0u\in W^{1,p}(\R^d)$ oraz
\[ \partial_j(E_0u)=E_0(\partial_j u)
\quad\text{w }\mathcal D'(\R^d),\qquad j=1,\ldots,d. \]
W konsekwencji normy $L^p$ funkcji i jej gradientu
są zachowane przez przedłużenie.
\end{wniosek}
\begin{proof}
Twierdzenie~\ref{thm:pde10-trace-zero} pozwala wybrać
$u_n\in C_c^\infty(\Omega)$ zbieżne do $u$ w $W^{1,p}(\Omega)$.
Przedłużenia przez zero $E_0u_n$ należą do
$C_c^\infty(\R^d)$ i spełniają
$\partial_j(E_0u_n)=E_0(\partial_j u_n)$ punktowo.
Z definicji przedłużenia
$\|E_0u_n-E_0u_m\|_{L^p(\R^d)}=
\|u_n-u_m\|_{L^p(\Omega)}$;
taka sama równość zachodzi dla gradientów.
Ciąg $(E_0u_n)$ jest zatem Cauchy'ego w $W^{1,p}(\R^d)$
i zbiega do pewnej funkcji $w$. Zbieżność w $L^p$ pokazuje,
że $w=E_0u$ prawie wszędzie. Zamkniętość pochodnej słabej
daje $\partial_jw=E_0(\partial_j u)$.
Równości norm wynikają z równości funkcji prawie wszędzie.
\end{proof}

\begin{uwaga}[Zależności twierdzeń i brzegi nieregularne]
Powyższy dowód świadomie korzysta z ustalonej wcześniej
równości $W_0^{1,p}(\Omega)=\ker\operatorname{Tr}$;
nie może więc zastąpić jej dowodu bez błędnego koła.
Na obszarze o nieregularnym brzegu nie należy
identyfikować obu przestrzeni bez dodatkowych hipotez.
\end{uwaga}

\section{Schemat związków i ograniczenia twierdzeń}
\label{sec:pde10-summary}
\begin{center}
\begin{tikzpicture}[node distance=8mm and 10mm,>=Stealth,
 every node/.style={font=\small,align=center}]
 \node[draw,rounded corners,fill=blue!5,text width=3.3cm] (a)
  {$W^{1,p}(\Omega)$\\$\Omega$ Lipschitzowska};
 \node[draw,rounded corners,fill=green!8,text width=3cm,
 below left=of a] (b) {$p<d$\\$L^{p^*}$ ciągle\\$L^q$, $q<p^*$ zwarto};
 \node[draw,rounded corners,fill=orange!10,text width=3cm,
 below=of a] (c) {$p=d$\\$L^q$ dla $q<\infty$\\zwarto dla $q<\infty$};
 \node[draw,rounded corners,fill=cyan!8,text width=3cm,
 below right=of a] (d) {$p>d$\\$C^{0,1-d/p}$ ciągle\\$C^{0,\beta}$ zwarto};
 \draw[->] (a)--(b);
 \draw[->] (a)--(c);
 \draw[->] (a)--(d);
\end{tikzpicture}

{\small\emph{Rysunek:} Zachowanie funkcji Sobolewa zależy od relacji
wykładnika całkowalności $p$ do wymiaru $d$. Zwarty charakter
zanurzeń zanika przy niektórych wykładnikach granicznych.}
\end{center}

Kluczowa różnica między zanurzeniem ciągłym a zwartym
polega na tym, że oszacowanie normy nie zapewnia istnienia
silnie zbieżnego podciągu. Przykład skupiania masy
pokazuje utratę zwartości przy $q=p^*$, a translacje
pokazują jej utratę na nieograniczonej przestrzeni.
Te dwa mechanizmy będą ważne w metodach wariacyjnych.


% --- Rozdział 11. Metody wariacyjne ---
